% META: {"id":"1SPE-VARALEA-EX-013","chapitre":"1SPE-VARIABLES-ALEATOIRES","type_objet":"exercice","capacites":["1SPE-VARIABLES-ALEATOIRES-C2"],"capacites_codes":["C2"],"parcours":1,"competences":["calculer","raisonner"],"duree_min":12,"mode_creation":"ex_nihilo","sources_inspiration":[],"fichier_tex":"chapitres/1SPE-VARIABLES-ALEATOIRES/exercices/1SPE-VARALEA-EX-013.tex","corrige_tex":"chapitres/1SPE-VARIABLES-ALEATOIRES/corriges/1SPE-VARALEA-CO-013.tex","status":"approved","origin":"generated","status_history":["generated","audited","accepted","approved"],"release_acceptance":"RELEASE_OWNER_BATCH_ACCEPTANCE","acceptance_closure_digest":"sha256:9b3ccf9a81c5520fbb7e03b7b2d3e7bf2057a02834b6d908bf3be5f49f3b553f","acceptance_receipt_digest":"sha256:4fad86f0282e0fa4e0419cca1ad6379ae7af5a280c04a4a90880f90a1c044242"}
% BEGIN-VERIFY
% from sympy import Rational, sqrt
% # X: loi donnee par tableau
% # x: -2, 0, 3 avec p: 1/4, 1/2, 1/4
% E = -2*Rational(1,4) + 0*Rational(1,2) + 3*Rational(1,4)
% assert E == Rational(1,4)
% E2 = 4*Rational(1,4) + 0*Rational(1,2) + 9*Rational(1,4)
% assert E2 == Rational(13,4)
% V = E2 - E**2
% assert V == Rational(13,4) - Rational(1,16)
% assert V == Rational(51,16)
% END-VERIFY
\begin{exercice}{1SPE-VARALEA-EX-013}{1}{12}
Soit $X$ une variable aléatoire de loi :
\[
\begin{array}{|c|c|c|c|}
\hline
x_i & -2 & 0 & 3 \\
\hline
P(X=x_i) & \frac{1}{4} & \frac{1}{2} & \frac{1}{4} \\
\hline
\end{array}
\]

\begin{enumerate}
  \item Calculer $E(X)$.
  \item Calculer $V(X)$ par la formule de König-Huygens.
  \item Calculer $\sigma(X)$.
\end{enumerate}
\end{exercice}
